\documentclass{beamer}
\usetheme{Madrid}

% pdflatex
\usepackage[utf8]{inputenc}
\usepackage[T2A]{fontenc}
\usepackage[russian,english]{babel}
\usepackage{cmap}
% \usepackage{lmodern}

\usepackage{graphicx}
\usepackage{xcolor}
\usepackage{tikz}
\usepackage{amsmath,amssymb,amsthm}
\usepackage{mathtools}
\usepackage{listings}
\usepackage{booktabs}
\usepackage{array}
\usetikzlibrary{arrows.meta,positioning,shapes.geometric,fit,backgrounds,calc,decorations.pathreplacing}

% Footline: show current section (left) and frame numbers (right)
\setbeamertemplate{footline}{%
  \leavevmode%
  \hbox{%
    \begin{beamercolorbox}[wd=.78\paperwidth,ht=2.5ex,dp=1ex,left]{author in head/foot}%
      \hspace*{1em}\usebeamerfont{footline}\insertsectionhead%
    \end{beamercolorbox}%
    \begin{beamercolorbox}[wd=.22\paperwidth,ht=2.5ex,dp=1ex,right]{date in head/foot}%
      \usebeamerfont{footline}\insertframenumber{} / \inserttotalframenumber\hspace*{1em}%
    \end{beamercolorbox}%
  }%
}

\newcommand{\parpause}[1]{\only<+->{#1\par}}
\newenvironment{definitionbox}[1][Определение]{%
  \begin{exampleblock}{#1}
}{%
  \end{exampleblock}
}

\AtBeginSection[]
{
  \begin{frame}
    \frametitle{Содержание}
    \tableofcontents[currentsection]
  \end{frame}
}

\title{Семинар 10. Еще немного о консенсусе: \\ византийские отказы и ABD}
\subtitle{Распределенные вычисления}
\author{Георгий Семенов}
\institute{Институт прикладных компьютерных наук \\ Университет ИТМО}
\date{весна 2026}

\begin{document}

\frame{\titlepage}

\begin{frame}
\frametitle{Организационные моменты}

\begin{itemize}
  \item Сегодня (23 апреля, 23:59) дедлайн по домашнему заданию 
  \item Домашнее задание 4 - Raft (дописать фрагменты кода, чтобы прошли тесты)
\end{itemize}

\end{frame}

%--------------------------------------------------
% \section{Paxos: Part-Time Parliament}

\begin{frame}
\frametitle{Какие алгоритмы консенсуса существуют?}


\begin{center}
\small
\begin{tabular}{>{\raggedright\arraybackslash}p{1.5cm} >{\raggedright\arraybackslash}p{3.0cm} c >{\raggedright\arraybackslash}p{4.0cm}}
\toprule
\textbf{Алгоритм} & \textbf{Автор(ы)} & \textbf{Год} & \textbf{Новшество} \\
\midrule
Paxos & Leslie Lamport & 1989/1998 & Безопасный консенсус в асинхронной модели с отказами (crash faults) \\
Raft & Diego Ongaro, John Ousterhout & 2014 & Упрощение понимания консенсуса через явное лидерство и репликацию лога \\
PBFT & Miguel Castro, Barbara Liskov & 1999 & Практический Byzantine Fault Tolerance для репликации состояния \\
ABD & Hagit Attiya, Amotz Bar-Noy, Danny Dolev & 1995 & Эмулируемый атомарный регистр на основе кворумов в message-passing системах \\
\bottomrule
\end{tabular}
\end{center}


\end{frame}


%--------------------------------------------------
\section{Византийский консенсус}

\begin{frame}
\frametitle{Проблема византийских генералов}

Слайды Мартина Клеппманна – 58-81 (проблема генералов + проблема византийских генералов)

\vspace{0.6cm}

\href{https://courses.cs.washington.edu/courses/cse452/20wi/slides/l15-bft.pdf}{Презентация про PBFT}


\end{frame}

%--------------------------------------------------
\section{ABD-протокол}

\begin{frame}
\frametitle{Линеаризуемость (Linearizability)}

\textbf{Определение:} существует глобальный порядок операций, согласованный с:
\begin{itemize}
  \item порядком операций внутри каждого потока (реального времени),
  \item ответом каждой операции.
\end{itemize}

\vspace{0.3em}
\small \textbf{Интуиция:} регистр «выглядит» как последовательное устройство, хотя на самом деле операции идут параллельно.

\vspace{0.6em}

Посмотреть по {\bfseries \href{https://eli.thegreenplace.net/2024/linearizability-in-distributed-systems/}{ссылке}}

\end{frame}

\begin{frame}
\frametitle{ABD: атомарные регистры без консенсуса}

\textbf{Задача:} реализовать линеаризуемый (атомарный) read/write регистр
в \emph{асинхронной} системе с crash-stop отказами.

\vspace{0.3em}
\begin{columns}[T]
\column{0.55\textwidth}
\begin{exampleblock}{Attiya, Bar-Noy, Dolev (1995)}
\small Регистр реплицируется на $n$ нодах. Операции Write и Read завершаются за 2 раунда. Корректно при $f < n/2$ отказов.
\end{exampleblock}

\vspace{0.4em}
\textbf{Чем отличается от консенсуса?}\\
\small Здесь не нужно навсегда согласиться на одном значении --- нужно лишь, чтобы чтения и записи \emph{выглядели} последовательными (linearizable).

\vspace{0.3em}
\textbf{Это \emph{слабее} консенсуса!}\\
\small Консенсус $\Rightarrow$ атомарный регистр,\\но регистр $\not\Rightarrow$ консенсус (в общем случае).

\column{0.42\textwidth}
\textbf{Зачем нужно:}
\begin{itemize}
  \item Репликация в БД
  \item Distributed key-value хранилище
  \item Distributed shared memory
\end{itemize}

\vspace{0.4em}
\begin{alertblock}{Связь с FLP}
\small ABD \textbf{решаем} в асинхронной модели! Потому что атомарный регистр $\neq$ консенсус.
\end{alertblock}
\end{columns}
\end{frame}

\begin{frame}
\frametitle{ABD: алгоритм записи}

\textbf{Модель:} $n$ реплик, каждая хранит пару $(\mathit{val},\ \mathit{tag})$, где $\mathit{tag} = (\mathit{ts},\ \mathit{writer\_id})$.

\vspace{0.3em}
\begin{block}{WRITE($v$) --- два раунда}
\begin{enumerate}
  \item \textbf{Read-phase:} спросить кворум о текущем максимальном теге $\Rightarrow$ получить $\mathit{max\_tag}$.
  \item \textbf{Write-phase:} записать $(v,\ \mathit{max\_tag}{+}1)$ на кворум $\Rightarrow$ ждать подтверждения.
\end{enumerate}
\end{block}

\vspace{0.1em}
\small Зачем Read-phase? Чтобы не затереть более свежее значение: нужно взять тег \emph{строго выше} всех существующих.

\vspace{0.3em}
\begin{center}
\begin{tikzpicture}[
  nd/.style={draw,rounded corners=3pt,minimum width=1.3cm,minimum height=0.5cm,
             align=center,font=\scriptsize},
  arr/.style={-{Stealth},thick,font=\scriptsize},
  scale=0.88, transform shape
]
  \node[nd,fill=blue!15] (cl) at (0,0) {Клиент};
  \node[nd,fill=green!20] (r1) at (-2.5,-1.6) {Реплика 1};
  \node[nd,fill=green!20] (r2) at (0,-1.6)    {Реплика 2};
  \node[nd,fill=green!20] (r3) at (2.5,-1.6)  {Реплика 3};

  % Phase 1
  \draw[arr,blue!60,bend right=8] (cl) to node[left,blue!60]{1. read\_tag} (r1);
  \draw[arr,blue!60]              (cl) -- node[left,blue!60]{} (r2);
  \draw[arr,blue!60,bend left=8]  (cl) to (r3);
  \draw[arr,blue!60,dashed,bend right=8] (r1) to node[right,blue!60]{\tiny max\_tag} (cl);
  \draw[arr,blue!60,dashed]              (r2) to (cl);

  % Phase 2
  \draw[arr,orange!80,bend right=10,transform canvas={yshift=-2pt}]
    (cl) to node[below left,orange!80,font=\tiny]{2. write(v, tag+1)} (r1);
  \draw[arr,orange!80,transform canvas={yshift=-2pt}]
    (cl) -- (r2);
  \draw[arr,orange!80,bend left=10,transform canvas={yshift=-2pt}]
    (cl) to (r3);

  \node[font=\tiny,orange!70] at (0,-2.55) {ждём подтверждения от кворума (2 из 3)};
\end{tikzpicture}
\end{center}
\end{frame}

\begin{frame}
\frametitle{ABD: алгоритм чтения и корректность}

\begin{block}{READ() --- два раунда}
\begin{enumerate}
  \item \textbf{Фаза 1:} спросить кворум $\Rightarrow$ собрать все $(v_i,\ \mathit{tag}_i)$ $\Rightarrow$ взять $(v^*,\ \mathit{tag}^*)$ с максимальным тегом.
  \item \textbf{Write-back:} записать $(v^*,\ \mathit{tag}^*)$ на кворум $\Rightarrow$ ждать подтверждения.
  \item Вернуть $v^*$.
\end{enumerate}
\end{block}

\vspace{0.2em}
\textbf{Зачем Write-back?}\\
\small Без него: читатель видит свежее значение у одной реплики, но другой читатель может уже не увидеть --- кворумы пересекаются, но «свежее» значение ещё не у всех. Write-back гарантирует: после чтения следующий читатель \emph{тоже} увидит это значение.

\vspace{0.3em}
\textbf{Почему корректно:} теги полностью упорядочены $\Rightarrow$ операция с большим тегом «позже». Кворумы пересекаются $\Rightarrow$ любой читатель видит ноду с последним записанным тегом.

\end{frame}

\begin{frame}
\frametitle{ABD: оптимизация для чтения}

\vspace{0.2em}
\textbf{Ограничения ABD:}
\begin{itemize}
  \item Базовая версия: \textbf{один писатель}; multi-writer требует тегов вида $(\mathit{ts},\ \mathit{writer\_id})$.
  \item Только read/write: \textbf{нельзя} реализовать compare-and-swap без консенсуса.
  \item 2 раунда на операцию --- дополнительная задержка.
\end{itemize}

\end{frame}

\begin{frame}
\frametitle{Итоги}

\begin{itemize}
    \item Выдано домашнее задание 4 (Raft)
    \item Византийский консенсус: проблема генералов, PBFT
    \item ABD-протокол: атомарные регистры без консенсуса, но с кворумами
\end{itemize}

\end{frame}

\end{document}